\unirsection{A fondo}

Los siguientes libros pueden servir de material de apoyo y para profundizar más en los contenidos de este tema.

\textbf{Nakahara, M. y Ohmi, T. Quantum Computing: From Linear Algebra to Physical Realizations. CRC Press, 2008.}

El capítulo 1 presenta las transformaciones lineales, la representación matricial y las matrices especiales (hermitianas, unitarias, normales) con una orientación directa hacia los circuitos cuánticos. Es la referencia más eficiente para conectar los resultados de este tema con sus aplicaciones en puertas cuánticas.

\textbf{Marinescu, D. y Marinescu, G. Lectures on Quantum Computing. University of Central Florida, 2003.}

Las secciones 3.1, 3.2 y 3.3 desarrollan los operadores lineales y su representación matricial en el contexto de la computación cuántica, con especial atención a los operadores hermíticos y unitarios. El tratamiento es más informal que el de este tema, lo que lo convierte en un complemento accesible.

\textbf{Axler, S. Linear Algebra Done Right (3.ª ed.). Springer, 2015.}

Los capítulos 3, 5 y 7 constituyen la referencia estándar moderna para transformaciones lineales, valores propios y el teorema espectral sobre espacios complejos. El enfoque, libre de determinantes en los primeros capítulos, es especialmente claro para comprender la estructura de los operadores hermitianos.

\textbf{Horn, R. A. y Johnson, C. R. Matrix Analysis (2.ª ed.). Cambridge University Press, 2013.}

Los capítulos 2 y 7 contienen un tratamiento riguroso y completo de la descomposición en valores singulares, las matrices unitarias y el teorema espectral complejo. Recomendado para quien desee profundizar en la SVD y sus propiedades de estabilidad numérica.